\documentclass[11pt]{article}
\usepackage[margin=0.85in]{geometry}
\usepackage{amsmath,amssymb,amsthm}
\usepackage[hidelinks]{hyperref}
\newcommand{\E}[1]{\mathrm{E}\left[#1\right]}
\newcommand{\PP}[1]{\mathrm{P}\left(#1\right)}
\newtheorem{theorem}{Theorem}
\newtheorem{lemma}{Lemma}
\newtheorem{corollary}{Corollary}
\setlength{\parskip}{3pt}
\setlength{\abovedisplayskip}{7pt plus 2pt minus 2pt}
\setlength{\belowdisplayskip}{7pt plus 2pt minus 2pt}
\begin{document}

\begin{center}
{\large\bfseries An Explicit Representation of the Cluster Duration}\\[3pt]
{Response counts, Dyck paths, and cyclic shifts}
\end{center}

Let $g(t)=\alpha e^{-\beta t}$ and $\rho=\alpha/\beta\in(0,1)$.
As in the main manuscript, $N$ includes the initiating point and
$\tau=\tau_{N-1}$ is the cluster duration. We prove an explicit
unconditional representation of $(N-1,\tau)$ as a compound Poisson
vector. The only new ingredient beyond the manuscript's conditional
representation is a weighted cyclic argument for response-count paths.

\paragraph{Response counts as Dyck paths.}
For a response-count vector $\mathbf k=(k_1,\ldots,k_n)$, write
$k_i^+=\sum_{j=1}^i k_j$ and $k_0^+=0$. Such a vector is the count
vector of a parking function of length $n$ exactly when
\begin{equation}\label{feasible}
k_n^+=n,\qquad k_i^+\ge i\quad(1\le i\le n).
\end{equation}
There are $n!/\prod_i k_i!$ parking functions with these counts, so
the uniform parking-function law in the main manuscript gives
\begin{equation}\label{countLaw}
\PP{\kappa=\mathbf k\mid N=n+1}
=\frac{n!}{(n+1)^{n-1}\prod_i k_i!}
\end{equation}
for every $\mathbf k$ satisfying \eqref{feasible}.

Draw $k_i$ up-steps followed by one down-step for each $i$. The height
after down-step $i$ is $k_i^+-i$, so \eqref{feasible} describes a Dyck
path. Its descendant schedule is $\Delta_i=k_i^+-i+1$. The conditional
independent-exponential representation in the main manuscript therefore
implies
\begin{equation}\label{conditionalLST}
\E{e^{-\theta\tau}\mid\kappa=\mathbf k,N=n+1}
=\prod_{i=1}^n
\frac{\beta(k_{i-1}^+-i+2)}{\beta(k_{i-1}^+-i+2)+\theta},
\qquad\theta\ge0.
\end{equation}
Indeed, the lists $(k_i^+-i)_{i=1}^n$ and
$(k_{i-1}^+-i+1)_{i=1}^n$ contain the same heights: they differ only
by replacing the final zero by the initial zero. Thus the rate
histogram $H$ is unchanged.

\paragraph{The weighted cyclic argument.}
For any nonnegative vector whose entries sum to $n$, write
\[
s_i=k_i^+-i,\qquad s_0=0,\qquad s_*=\min_{0\le r<n}s_r.
\]
Thus $s_i$ is the path height, and a vector is feasible exactly when
$s_i\ge0$, in which case $s_*=0$. For a feasible $\mathbf k$, define
\begin{equation}\label{weight}
w(\mathbf k)=z^n\prod_{i=1}^n
\frac{e^{-\rho}\rho^{k_i}}{k_i!}
\frac{\beta(k_{i-1}^+-i+2)}{\beta(k_{i-1}^+-i+2)+\theta},
\qquad 0\le z\le1,
\end{equation}
and let $h(\mathbf k)$ be the number of returns of its Dyck path to
zero, including the final return.

\begin{lemma}[Weighted cyclic counting]\label{cyclicLemma}
For every $n\ge1$,
\begin{equation}\label{cyclicEq}
\sum_{\substack{k_n^+=n\\k_i^+\ge i}}
\frac{w(\mathbf k)}{h(\mathbf k)}
=\frac{z^n}{n}
\sum_{\substack{k_1+\cdots+k_n=n\\k_i\ge0}}
\prod_{i=1}^n\frac{e^{-\rho}\rho^{k_i}}{k_i!}
\frac{\beta(1+s_{i-1}-s_*)}
{\beta(1+s_{i-1}-s_*)+\theta}.
\end{equation}
\end{lemma}
\begin{proof}
Fix a nonnegative vector whose entries sum to $n$. Its increments
$k_i-1$ sum to zero. A cyclic shift beginning after position $r$ is a
Dyck path exactly when $s_r=s_*$. If this minimum occurs $h$ times, exactly $h$
of the $n$ starting positions give Dyck paths, each with $h$ returns
to zero.

Cyclic shifts preserve $\prod_i(e^{-\rho}\rho^{k_i}/k_i!)$.
They also cyclically permute the path heights after their common
minimum has been subtracted, and hence preserve the product of the
Laplace-transform factors in \eqref{cyclicEq}. Therefore averaging
over all $n$ starting positions contributes $(h/n)(1/h)=1/n$ for
each unrestricted vector, proving \eqref{cyclicEq}. Starting positions
are counted separately even when two shifts are identical, so periodic
vectors cause no exception.
\end{proof}

\begin{theorem}[Joint compound Poisson representation]\label{jointDecomp}
Let $B\sim\mathsf{Borel}(\rho)$. Conditional on $B=n$, let
\[
(K_1,\ldots,K_n)\sim\mathsf{Multinomial}(n;1/n,\ldots,1/n),
\qquad K_i^+=\sum_{j=1}^iK_j,\quad K_0^+=0,
\]
and let $E_1,E_2,\ldots$ be independent $\mathsf{Exp}(1)$ variables,
independent of $B$ and the counts. Define
\begin{equation}\label{markDef}
J=\frac1\beta\sum_{i=1}^B
\frac{E_i}{K_{i-1}^+-i+2-\min_{0\le r<B}(K_r^+-r)},
\end{equation}
where an empty sum is zero. If $M\sim\mathsf{Poisson}(\rho)$ and
$(B_\ell,J_\ell)$ are independent copies of $(B,J)$, independent of
$M$, then
\begin{equation}\label{jointDecompEq}
(N-1,\tau)\stackrel{\mathsf D}{=}
\sum_{\ell=1}^M(B_\ell,J_\ell).
\end{equation}
The sum is $(0,0)$ when $M=0$.
\end{theorem}

Here is how to read the four distributions in the theorem. The outer
Poisson variable $M$ comes from taking the logarithm of the Dyck-path
generating function: the logarithm isolates the cyclic pieces, and
exponentiating their transform produces a Poisson sum. A piece of
length $n$ receives the coefficient
$n^{-1}\PP{\mathsf{Poisson}(n\rho)=n}$, which is exactly
$\rho\PP{B=n}$; this is why its size is Borel. Before the cyclic shift,
the $n$ entries are independent $\mathsf{Poisson}(\rho)$ response
counts conditioned to sum to $n$. Equal Poisson means make their
conditional law $\mathsf{Multinomial}(n;1/n,\ldots,1/n)$, which
explains the count mixture. Finally, after shifting the count path at
its minimum, its positive heights are the numbers of pending responses.
Exponential memorylessness makes the waiting time at height $j$ an
$\mathsf{Exp}(\beta j)$ variable. Thus $J$ is the corresponding sum of
independent exponential waiting times. In this order, the construction is
\[
\begin{gathered}
\text{Poisson number of pieces}
\;\longrightarrow\;\text{Borel size of each piece}\\
\longrightarrow\;\text{multinomial response counts}
\;\longrightarrow\;\text{exponential waiting times}.
\end{gathered}
\]
This description explains the mixture; it does not identify the pieces
with physical descendant subclusters.

\begin{proof}
Let $F(z,\theta)=\E{z^{N-1}e^{-\theta\tau}}$. Multiplying
\eqref{countLaw} by the Borel probability of $N=n+1$ and using
\eqref{conditionalLST} gives
\[
F(z,\theta)=e^{-\rho}
\left(1+\sum_{n\ge1}\sum_{\substack{k_n^+=n\\k_i^+\ge i}}
w(\mathbf k)\right).
\]
Every Dyck path splits uniquely at its returns to zero into an ordered
sequence of nonempty paths with no intermediate return. The weight in
\eqref{weight} factors over this sequence because each piece restarts
at height zero. If $R$ is the total weight of one such piece, then
\[
F(z,\theta)=\frac{e^{-\rho}}{1-R},
\qquad
\log F(z,\theta)=-\rho+\sum_{h\ge1}\frac{R^h}{h}.
\]
At $z=1$ and $\theta=0$, $R=1-e^{-\rho}<1$; elsewhere in the stated
domain the nonnegative weights are no larger. Hence rearrangement is
valid. The term $R^h/h$ assigns $w(\mathbf k)/h(\mathbf k)$ to each
path. Lemma~\ref{cyclicLemma} therefore removes the path restrictions:
\begin{align}\label{logF}
\log F(z,\theta)
&=-\rho+\sum_{n\ge1}\frac{z^n}{n}
\sum_{\substack{k_1+\cdots+k_n=n\\k_i\ge0}}
\prod_{i=1}^n\frac{e^{-\rho}\rho^{k_i}}{k_i!}\nonumber\\[-2pt]
&\hspace{28mm}\times
\prod_{i=1}^n
\frac{\beta(1+s_{i-1}-s_*)}
{\beta(1+s_{i-1}-s_*)+\theta}.
\end{align}
The first product is the joint probability of $n$ independent
$\mathsf{Poisson}(\rho)$ variables. Conditional on their sum being
$n$, they have the multinomial law in the theorem; the second product
is the conditional Laplace transform of $J$. Finally,
\[
\frac1n\PP{\mathsf{Poisson}(n\rho)=n}
=\frac{e^{-n\rho}(n\rho)^n}{n\,n!}
=\rho\,\PP{B=n}.
\]
Substitution in \eqref{logF} yields
\[
F(z,\theta)=
\exp\left\{\rho\left(\E{z^B e^{-\theta J}}-1\right)\right\},
\]
the joint transform of the Poisson sum in \eqref{jointDecompEq}.
Uniqueness of the transform proves the result.
\end{proof}

\begin{corollary}\label{durationDecomp}
The duration is compound Poisson:
\[
\tau\stackrel{\mathsf D}{=}\sum_{\ell=1}^M J_\ell,
\qquad M\sim\mathsf{Poisson}(\rho),
\]
with the explicit mark law \eqref{markDef}. In particular,
$\PP{\tau=0}=e^{-\rho}$.
\end{corollary}
\begin{proof}
Take the second coordinate of \eqref{jointDecompEq}. Since every
denominator in \eqref{markDef} is a positive integer, $J>0$ almost
surely, so the sum is zero exactly when $M=0$.
\end{proof}

Corollary~\ref{durationDecomp} is an unconditional generative
description of the duration. First draw $M$, the number of independent
contributions. For each contribution, draw its number $B$ of exponential
terms; the Borel law appears because the coefficient of a cyclic piece
of length $n$ is $\rho\PP{B=n}$. Given that length, distribute $n$
responses among its $n$ positions according to the multinomial law.
Shift the resulting count path at a minimum so that all pending-response
counts are positive, and then draw one exponential waiting time at each
height. Their sum is $J$. Adding the $M$ independent values of $J$
produces $\tau$ in distribution.

This makes several features immediate. The duration is infinitely
divisible and has an atom of size $e^{-\rho}$ at zero, corresponding to
no contribution. Conditional on the multinomial counts, repeated visits
to height $j$ may be grouped into a gamma variable with shape $H_j$ and
rate $\beta j$, exactly as in the descendant-schedule representation in
the main manuscript. The Borel variable determines how many waiting
times occur, the multinomial vector determines their rates, and the
exponential variables determine their realized lengths. The construction
is an exact sampling recipe even though it does not give an elementary
closed-form density or CDF.

\paragraph{Recovery of the conditional mean.}
We next verify explicitly that the representation reproduces the mean
identities in the main manuscript. Throughout this paragraph condition
on $N=n+1$. Differentiating \eqref{conditionalLST} at $\theta=0$ gives
\begin{equation}\label{conditionalMean}
\beta\E{\tau\mid N=n+1}
=\sum_{m=0}^{n-1}\E{\frac1{\kappa_{m+1}^+-m}\,\bigg|\,N=n+1}.
\end{equation}
The generalized parking-function count used in the manuscript gives,
for $c=\kappa_m^+$,
\begin{align}\label{nextCount}
&\PP{\kappa_{m+1}=k\mid
\kappa_1,\ldots,\kappa_m,N=n+1}\nonumber\\[-2pt]
&\quad=
\frac{c+k-m}{c-m+1}\binom{n-c}{k}
\frac{(n-m)^{n-c-k-1}}{(n-m+1)^{n-c-1}}.
\end{align}
Divide \eqref{nextCount} by $c+k-m$ and sum over $k$. The linear
factor cancels, and the binomial theorem gives
\begin{align}\label{innerMean}
&\E{\frac1{\kappa_{m+1}^+-m}\,\bigg|
\kappa_1,\ldots,\kappa_m,N=n+1}\nonumber\\[-2pt]
&\quad=
\frac{n-m+1}{(n-m)(\kappa_m^+-m+1)}
-\left(\frac{n-m}{n-m+1}\right)^{n-m-1}
\mathbf1\{\kappa_m^+=m\}.
\end{align}
The subtracted term is exactly the inadmissible $k=0$ term when
$c=m$: with no response pending, the next count must be positive.

For completeness, the event $\kappa_m^+=m$ splits the parking function
at $m$. Choose the $m$ coordinates occupied by values at most $m$,
then choose an ordinary parking function on each side. Since there are
$(r+1)^{r-1}$ parking functions of length $r$,
\begin{equation}\label{returnProbability}
\PP{\kappa_m^+=m\mid N=n+1}
=\binom nm\frac{(m+1)^{m-1}(n-m+1)^{n-m-1}}
{(n+1)^{n-1}}.
\end{equation}
Taking expectations in \eqref{innerMean} and substituting
\eqref{returnProbability} recovers the manuscript's recurrence
\begin{align}\label{meanRecurrence}
\E{\frac1{\kappa_{m+1}^+-m}\,\bigg|N=n+1}
&=\frac{n-m+1}{n-m}
\E{\frac1{\kappa_m^+-m+1}\,\bigg|N=n+1}\nonumber\\[-2pt]
&\quad-
\binom nm\frac{(m+1)^{m-1}(n-m)^{n-m-1}}
{(n+1)^{n-1}}.
\end{align}
Here $\kappa_0^+=0$. Multiplying by $n-m$ and iterating from $m=0$
shows explicitly that
\begin{align}\label{meanRecurrenceSolution}
\E{\frac1{\kappa_{m+1}^+-m}\,\bigg|N=n+1}
&=\frac{n+1}{n-m}\nonumber\\[-2pt]
&\quad-\frac1{n-m}\sum_{i=0}^m\binom ni
\left(\frac{i+1}{n+1}\right)^{i-1}
\left(\frac{n-i}{n+1}\right)^{n-i}.
\end{align}
Indeed, substituting $m-1$ into \eqref{meanRecurrenceSolution} and
then applying \eqref{meanRecurrence} adds precisely the $i=m$ term.
Equations \eqref{conditionalMean} and \eqref{meanRecurrenceSolution}
are the conditional mean formula in the main manuscript.

\paragraph{Recovery of the Borel mean identity.}
Now remove the conditioning. Since
\[
\PP{N=n+1}=\frac{e^{-\rho(n+1)}(\rho(n+1))^n}{(n+1)!},
\]
substituting \eqref{meanRecurrenceSolution} into
\eqref{conditionalMean}, multiplying by this probability, and summing
over $n\ge1$ produces two terms. The first is
\begin{align}
\sum_{n\ge1}(n+1)\mathsf H_n
\frac{e^{-\rho(n+1)}(\rho(n+1))^n}{(n+1)!}
=\E{B\mathsf H_{B-1}\mathbf1\{B\ge2\}},
\label{firstMeanTerm}
\end{align}
where $B\sim\mathsf{Borel}(\rho)$ and
$\mathsf H_j=\sum_{r=1}^j1/r$. For the second term, interchange the
$m$- and $i$-sums, so that
$\sum_{m=i}^{n-1}1/(n-m)=\mathsf H_{n-i}$, and then set
$j=i+1$ and $\ell=n-i$. It factors as
\begin{align}
&\sum_{i\ge0}
\rho\frac{e^{-\rho(i+1)}(\rho(i+1))^i}{(i+1)!}
\sum_{\ell\ge1}\ell\mathsf H_\ell
\frac{e^{-\rho\ell}(\rho\ell)^{\ell-1}}{\ell!}\nonumber\\[-2pt]
&\qquad=\rho\,\E{B\mathsf H_B},
\label{secondMeanTerm}
\end{align}
because the first sum after the factor $\rho$ is the total Borel
probability. Therefore \eqref{firstMeanTerm}--\eqref{secondMeanTerm}
give
\begin{align*}
\beta\E{\tau}
&=\E{B\mathsf H_{B-1}\mathbf1\{B\ge2\}}
-\rho\E{B\mathsf H_B}\\
&=(1-\rho)\E{B\mathsf H_B}-1.
\end{align*}
The last equality uses $B\mathsf H_{B-1}=B\mathsf H_B-1$ for
$B\ge2$ and checks $B=1$ separately. Finally, Theorem~\ref{jointDecomp}
also gives $\E{\tau}=\rho\E{J}$, so the new decomposition recovers
the equivalent mark identity
\begin{equation}\label{markMean}
\beta\rho\E{J}=(1-\rho)\E{B\mathsf H_B}-1.
\end{equation}

The integral form follows from the same identity without another
combinatorial calculation. Let $b(x)=\E{x^B}$. The Borel generating
function satisfies $b(x)=x\exp\{\rho(b(x)-1)\}$, and differentiation
gives
\[
xb'(x)=\frac{b(x)}{1-\rho b(x)},
\qquad \E{B}=b'(1)=\frac1{1-\rho}.
\]
Using $\mathsf H_n=\int_0^1(1-x^n)/(1-x)\,\mathrm dx$ in the Borel
mean identity gives
\begin{align*}
\beta\E{\tau}
&=(1-\rho)\int_0^1
\frac{\E{B}-\E{Bx^B}}{1-x}\,\mathrm dx-1\\
&=\int_0^1\frac{1-(1-\rho)xb'(x)}{1-x}\,\mathrm dx-1.
\end{align*}
Set $u=\rho b(x)$. The defining equation for $b$ implies
\[
x=\frac{ue^{-u}}{\rho e^{-\rho}},
\qquad
\mathrm dx=\frac{(1-u)e^{-u}}{\rho e^{-\rho}}\,\mathrm du,
\]
and $x=0,1$ correspond to $u=0,\rho$. Direct substitution yields
\begin{align*}
\beta\E{\tau}
&=\frac1\rho\int_0^\rho
\frac{(\rho-u)e^{-u}}{\rho e^{-\rho}-ue^{-u}}\,\mathrm du-1\\
&=\int_0^\rho
\frac{e^{-u}-e^{-\rho}}{\rho e^{-\rho}-ue^{-u}}\,\mathrm du.
\end{align*}
At $u=\rho$ the integrand is understood by continuity. This recovers
the integral identity in the main manuscript as well as its Borel form.
Setting $\theta=0$ in the joint transform also recovers
$\E{z^N}=\E{z^B}$, while
$J\le\beta^{-1}\sum_{i=1}^B E_i$ confirms that the displayed mean is
finite.

\end{document}
